% titel: Blatt 10, A4: Gleichgewichtslage vorhersagen (Säure/Basen)
% skript: Mortimer Kap. 17 (Säurestärke, Gleichgewicht)

% Teil I: Werkzeuge, kein Lösungsweg.

Bei einer Protonen-Transfer-Reaktion liegt das Gleichgewicht auf der Seite
der \emph{schwächeren} Säure und schwächeren Base (Reaktion bildet das
stabilere Paar).

\begin{werkzeugkasten}[Gleichgewichtslage]
Im Säure-Base-Gleichgewicht wandert das Proton zur stärkeren Base; es
entsteht die schwächere Säure.
\[
  \underbrace{\text{Säure$_1$} + \text{Base$_2$}}_{\text{rechts, wenn
  Säure$_1$ stark}} \rightleftharpoons
  \text{Base$_1$} + \underbrace{\text{Säure$_2$}} .
\]
Gleichgewicht auf der Seite der schwächeren Säure/Base.
Hilfreich: $pK_a$-Werte bzw.\ qualitative Stärke (Strukturregeln des
vorherigen Aufgabentyps) der beteiligten Säuren vergleichen.
\textit{Skript:} Mortimer Kap. 17 (Säure-Base-Gleichgewichte).
\end{werkzeugkasten}

\begin{achtung}[Darauf achten]
\begin{itemize}
  \item Konjugierte Paare bilden: HSO$_4^-$/$SO_4^{2-}$,
        HCO$_3^-$/CO$_3^{2-}$, NH$_4^+$/NH$_3$, H$_2$PO$_4^-$/$HPO_4^{2-}$.
  \item Frage ist \emph{nicht} nach $K$-Wert, sondern nach der Richtung.
  \item Wohin ein Proton wandert (stärkste Base zieht es an), bestimmt die
        bevorzugte Seite.
  \item Wasser (H$_2$O$\rightleftharpoons$OH$^-$, $pK_a\approx14$) als
        schwache Säure einordnen.
\end{itemize}
\end{achtung}
